% Options for packages loaded elsewhere
\PassOptionsToPackage{unicode}{hyperref}
\PassOptionsToPackage{hyphens}{url}
\documentclass[
  ukrainian,
]{article}
\usepackage{xcolor}
\usepackage[margin=2.5cm]{geometry}
\usepackage{amsmath,amssymb}
\setcounter{secnumdepth}{-\maxdimen} % remove section numbering
\usepackage{iftex}
\ifPDFTeX
  \usepackage[T1]{fontenc}
  \usepackage[utf8]{inputenc}
  \usepackage{textcomp} % provide euro and other symbols
\else % if luatex or xetex
  \usepackage{unicode-math} % this also loads fontspec
  \defaultfontfeatures{Scale=MatchLowercase}
  \defaultfontfeatures[\rmfamily]{Ligatures=TeX,Scale=1}
\fi
\usepackage{lmodern}
\ifPDFTeX\else
\setmainfont{DejaVu Serif}
\setsansfont{DejaVu Sans}
\setmonofont{DejaVu Sans Mono}
\setmathfont{STIX Two Math}
\fi
\ifPDFTeX\else
  % xetex/luatex font selection
\fi
% Use upquote if available, for straight quotes in verbatim environments
\IfFileExists{upquote.sty}{\usepackage{upquote}}{}
\IfFileExists{microtype.sty}{% use microtype if available
  \usepackage[]{microtype}
  \UseMicrotypeSet[protrusion]{basicmath} % disable protrusion for tt fonts
}{}
\makeatletter
\@ifundefined{KOMAClassName}{% if non-KOMA class
  \IfFileExists{parskip.sty}{%
    \usepackage{parskip}
  }{% else
    \setlength{\parindent}{0pt}
    \setlength{\parskip}{6pt plus 2pt minus 1pt}}
}{% if KOMA class
  \KOMAoptions{parskip=half}}
\makeatother
\ifLuaTeX
\usepackage[bidi=basic,shorthands=off,provide=*]{babel}
\else
\usepackage[bidi=default,shorthands=off,provide=*]{babel}
\fi
\ifLuaTeX
  \usepackage{selnolig} % disable illegal ligatures
\fi
\setlength{\emergencystretch}{3em} % prevent overfull lines
\providecommand{\tightlist}{%
  \setlength{\itemsep}{0pt}\setlength{\parskip}{0pt}}
\usepackage{bookmark}
\IfFileExists{xurl.sty}{\usepackage{xurl}}{} % add URL line breaks if available
\urlstyle{same}
\hypersetup{
  pdflang={uk},
  hidelinks,
  pdfcreator={LaTeX via pandoc}}

\author{}
\date{}

\begin{document}

\textbf{РИС.} 42.3. Блок-схема розрахунку інтегрального
\texttt{\textbackslash{}int\_\{a\}\^{}\{b\}\ f(x)\ dx} методом
Монте-Карло (100 кидків);
\texttt{M\ =\ \textbackslash{}max\_\{a\}\ f(x)} .

\section{Випадкові
числа}\label{ux432ux438ux43fux430ux434ux43aux43eux432ux456-ux447ux438ux441ux43bux430}

Перш ніж перейти до рівняння Монте-Карло з частковими похідними,
зупинимося на важливому питанні випадкових чисел. Щойно, обчислюючи
інтеграл, нам потрібно було побудувати послідовність випадкових точок
\texttt{P\_i\ =\ (x\_i,\ y\_i)} всередині прямокутника R. Іншими
словами, x-координата точки P має бути випадковим числом з інтервалу
{[}a, b{]}, а \texttt{y\_i} від інтервалу {[}0, M{]}. Щоб знайти
випадкові числа з цих інтервалів, звертаємося до послідовності
випадкових чисел, рівномірно розподілених по сегменту {[}0, 1{]}. Тоді
числа, рівномірно розподілені на сегменті {[}a, b{]}, можна обчислити за
формулою

(початкове зміщення струни), \texttt{u\_t}(x, 0) = \texttt{g(x)}
(початкова швидкість струни).

Таким чином, хвильове рівняння відрізняється від рівняння
теплопровідності, для якого потрібно встановити лише одну початкову
умову.

\begin{enumerate}
\def\labelenumi{\arabic{enumi}.}
\setcounter{enumi}{3}
\tightlist
\item
  За допомогою хвильового рівняння можна описати розподіл електричного
  струму в дроті. З законів Кірхгофа отримуємо таку систему двох рівнянь
  у похідних похідних першого порядку:
\end{enumerate}

(16.3)

Тепер, звичайно, виникає питання: як отримати послідовність випадкових
чисел,
\texttt{\textbackslash{}\{r\_i:\ i=1,2,\textbackslash{}ldots\textbackslash{}\}}
рівномірно розподілених на {[}0,1{]}? Найпоширенішим методом сьогодні є
метод дедукцій (метод порівнянь). Використовуючи цей метод, отримуємо
послідовність випадкових цілих чисел (наприклад, 2120, 1401, 177,
3013,\ldots), потім ліворуч від кожного числа розміщується десяткова
крапка, так що всі числа знаходяться між нулем і одним
\texttt{(0,212;\ 0,1401;\ 0,0177;\ 0,3013;\ \textbackslash{}ldots)} .

Отже, щоб отримати послідовність випадкових цілих чисел (між 0 і P), ми
використаємо метод віднімання.

\section{Отримання випадкових чисел за допомогою методу
відрахування.}\label{ux43eux442ux440ux438ux43cux430ux43dux43dux44f-ux432ux438ux43fux430ux434ux43aux43eux432ux438ux445-ux447ux438ux441ux435ux43b-ux437ux430-ux434ux43eux43fux43eux43cux43eux433ux43eux44e-ux43cux435ux442ux43eux434ux443-ux432ux456ux434ux440ux430ux445ux443ux432ux430ux43dux43dux44f.}

\begin{enumerate}
\def\labelenumi{\arabic{enumi}.}
\item
  Як перше випадкове ціле число, обирайте будь-яке число, що знаходиться
  між 0 і P (число P обирається наперед).
\item
  Помножте це випадкове число на певне фіксоване число.
\end{enumerate}

Попередньо обраний множник M.

\begin{enumerate}
\def\labelenumi{\arabic{enumi}.}
\setcounter{enumi}{2}
\item
  Додайте до добутку деяке фіксоване ціле число \emph{K}, обране
  заздалегідь.
\item
  Поділити отриману суму на P і залишок
\end{enumerate}

Виберіть нове випадкове число.

Метод виведення можна записати у формі формули
\texttt{r\_\{i+1\}\ \textbackslash{}equiv\ (Mr\_i\ +\ K)\ \textbackslash{}bmod\ P}
, що означає, що для отримання наступного \texttt{r\_\{i+1\}} потрібно
взяти попередній \texttt{r\_i} , помножити його на M, додати результат
на K, поділити суму на P і взяти решту поділу. Наприклад, якщо P = 100,
M = 37, K = 16, \texttt{r\_0\ =\ 15} , то отримується наступна
послідовність випадкових чисел:

где

ж --- координата вдоль провода,

\emph{t} --- время,

\begin{verbatim}
i(x, t) — распределение тока вдоль провода,  v\left(x,\ t\right)  — распределение потенциала вдоль провода,
\end{verbatim}

С -- емкость провода на единицу длины,

G --- утечка на единицу длины,

Р --- сопротивление на единицу длины,

-- индуктивность провода на единицу длины.

Рівняння (16.3) зазвичай називають системою телеграфних рівнянь. У такій
формі вони зазвичай не використовуються. Щоб перетворити цю систему,
диференціюйте перше рівняння на x, друге на t, помножте обидві частини
на C і відніміть від першого. У результаті ми отримуємо

\begin{verbatim}
Mr_1 + K = 2643
\end{verbatim}

Воспользуемся теперь вторым уравнением из (16.3)

\begin{verbatim}
r = 71
\end{verbatim}

окончательно получаем

(16.4)

Y

Це рівняння для струму, відоме як телеграфне рівняння, є гіперболічним
рівнянням другого порядку (якщо, звісно, C і L не дорівнюють нулю,
інакше рівняння стає параболічним).

Напряжение описывается точно таким же уравнением:

(16.5)

\begin{verbatim}
I = \int_{0}^{1} \int_{0}^{1} \int_{0}^{1} \int_{0}^{1} e^{-(x^{2}+y^{2}+z^{2}+\omega^{2})} dx dy dz dw,
\end{verbatim}

Если G = R = 0, то уравнения (16.4)---(16.5) упрощаются:

(16.6)

\begin{verbatim}
I = \int_{0}^{1} e^{\sin x} dx.
\end{verbatim}

\section{\texorpdfstring{ЗАВДАННЯ - 1. Отримайте рівняння (16.5) для
\emph{and} з системи з двох рівнянь першого порядку
(16.3)}{ЗАВДАННЯ - 1. Отримайте рівняння (16.5) для and з системи з двох рівнянь першого порядку (16.3)}}\label{ux437ux430ux432ux434ux430ux43dux43dux44f---1.-ux43eux442ux440ux438ux43cux430ux439ux442ux435-ux440ux456ux432ux43dux44fux43dux43dux44f-16.5-ux434ux43bux44f-and-ux437-ux441ux438ux441ux442ux435ux43cux438-ux437-ux434ux432ux43eux445-ux440ux456ux432ux43dux44fux43dux44c-ux43fux435ux440ux448ux43eux433ux43e-ux43fux43eux440ux44fux434ux43aux443-16.3}

\begin{itemize}
\tightlist
\item
  \begin{enumerate}
  \def\labelenumi{\arabic{enumi}.}
  \setcounter{enumi}{1}
  \tightlist
  \item
    Знаючи фізичне значення кожного члена хвильового рівняння, що ви
    можете сказати про поведінку в момент розв'язання наступної задачі:
  \end{enumerate}
\end{itemize}

(YYII)

\begin{verbatim}
r_{i+1} \equiv (3r_i + 4) \mod 7, \quad r_0 = 0.
,  0 < x < 1  ,  0 < t < \infty  ,
\end{verbatim}

(ГУ)

\$\$
\texttt{I\ =\ \textbackslash{}int\_\{0\}\^{}\{1\}\ \textbackslash{}int\_\{0\}\^{}\{1\}\ \textbackslash{}int\_\{0\}\^{}\{1\}\ e\^{}\{-(x\^{}\{2\}+y\^{}\{2\}+z\^{}\{2\})\}\ dx\ dy\ dz.}
u(x,y) = g(x,y) =
\texttt{\textbackslash{}begin\{cases\}\ u(0,\ t)\ =\ 0,\ \&\ 0\ \textless{}\ t\ \textless{}\ \textbackslash{}infty,\ \textbackslash{}\textbackslash{}\ u(1,\ t)\ =\ 0,\ \&\ 0\ \textless{}\ t\ \textless{}\ \textbackslash{}infty,\ \textbackslash{}end\{cases\}}
Після того, як цю конкретну задачу буде розв'язано методом Монте-Карло,
ми розглянемо, як розв'язувати більш загальні задачі за допомогою цього
методу.

Щоб проілюструвати метод Монте-Карло, розглянемо гру під назвою
«Блукаючий п'яниця». Щоб грати в цю гру, потрібна таблиця з сіткою
ліній, як показано на рис. 43.1.

Перейдем к описанию правил игры.

\end{document}
